\chapter{Le volet des modèles de langue}
\label{chap:07-modeles-de-langue}

Les chapitres précédents traitent de programmes qui apprennent à agir dans un monde simulé. Ce
chapitre décrit un volet plus petit, consacré aux modèles d'IA générative qui écrivent du texte. Il
présente le vocabulaire retenu, l'algorithme tel que nous l'avons écrit, le protocole, les essais et
leurs résultats, puis le piège de mesure que ces essais ont révélé. Le protocole complet se trouve
dans \file{sim/llm/LLM\_PROTOCOL.md}. Ses vingt écarts sont résumés au \cref{chap:09-ecarts}. Les écarts~18
et~19, écrits le 21 novembre 2026 après tous les résultats, déclarent des défauts trouvés par une relecture
technique et tirent des journaux des mesures que la première version de ce chapitre laissait de côté. Ils
changent la lecture du volet sans changer aucun chiffre publié.

\section{Le vocabulaire}
\label{sec:llm-vocabulaire}

Le vocabulaire du RL a gagné les modèles d'IA générative. On emprunte au RL certains de ses
algorithmes pour affiner les réponses d'un modèle de langue, sans agent qui agit dans un monde. Un modèle de langue est un modèle IA qui produit un texte jeton par
jeton. Un jeton est un morceau de mot. Pour un problème de mathématiques, un programme compare la
réponse finale à la bonne réponse et la récompense vaut 1 quand elles concordent. La littérature
parle de RL à récompense vérifiable (RLVR, \emph{reinforcement learning with verifiable rewards}).
Nous gardons ce sigle quand nous citons un article qui l'emploie, sans en reprendre l'idée qu'il
s'agit du même problème que celui d'un agent dans un monde.

Trois différences séparent ce cadre du RL classique des \cref{chap:05-experiences-methode,chap:06-experiences-resultats} :
\begin{itemize}
  \item un \enquote{épisode} se réduit à une seule réponse : aucun état n'évolue sous l'effet des
        actions, aucune récompense n'est différée sur plusieurs décisions ;
  \item la politique de départ est un modèle déjà entraîné sur d'immenses corpus, qui sait souvent
        répondre avant tout affinage ;
  \item la récompense ne dépend que du texte final.
\end{itemize}
Nous n'appelons jamais RL l'affinage par retours humains (RLHF,
\emph{reinforcement learning from human feedback}, \cite{ouyang2022instructgpt}) : nous le
désignons par son sigle, comme une technique d'affinage parmi d'autres.

\section{Pourquoi ce volet}
\label{sec:llm-pourquoi}

En 2025, \textcite{shao2025spurious} ont montré qu'une récompense tirée au sort, sans regarder si la
réponse est juste, améliore les modèles Qwen en mathématiques, mais pas les modèles Llama ou OLMo.
Sur le modèle Qwen2.5-Math-7B et le banc d'essai MATH-500, la vraie récompense fait gagner
29,1 points de pourcentage. Une récompense qui félicite une réponse fausse fixée d'avance en fait
gagner 24,1, une récompense tirée au hasard 21,4. Les auteurs l'expliquent par un comportement que
l'entraînement amplifie sans le créer : le modèle écrit son raisonnement sous forme de code, sans
l'exécuter, dans 65~\% des réponses avant l'entraînement et dans plus de 90~\% après un entraînement
à récompense trompeuse. Ils recommandent de valider une méthode sur plusieurs familles de modèles.

Ce résultat prolonge deux failles de 2018 : le meilleur algorithme change avec la tâche (F13) et les
bancs d'essai trop simples se laissent apprendre (F16). Le sondage du \cref{chap:04-sondage-resultats}
mesure le même risque dans les articles : seize articles sur 27 du premier tirage de 2026 qui
appliquent ce qu'ils appellent RL à de grands modèles de langue n'essaient qu'une seule famille de
modèles (critère X6). Notre volet regarde le phénomène de près, à petite échelle, avec deux familles
de modèles. Le protocole annonce d'avance sa portée : une illustration sur un ordinateur portable. Elle ne
réplique pas \textcite{shao2025spurious}.

\section{L'algorithme GRPO, tel que nous l'avons écrit}
\label{sec:grpo}

GRPO (\emph{Group Relative Policy Optimization}) vient de \textcite{shao2024deepseekmath}. Il
reprend la fonction objectif coupée de PPO \parencite{schulman2017ppo} et remplace le réseau
critique, qui estime la valeur d'un état, par une normalisation au sein d'un groupe de réponses au
même problème. Nous l'avons écrit à la main, en PyTorch, dans un seul fichier, \file{sim/llm/grpo.py}, pour
contrôler chaque étape. Le fichier comptait 159 lignes pour les essais décrits ici. Il en compte 266 depuis
l'ajout de la relance de l'écart~20, derrière l'option \texttt{-{}-c7}. Sans cette option, il refait les
anciens essais à l'identique. Aucune bibliothèque d'affinage n'intervient :
le fichier n'importe que PyTorch et la bibliothèque Transformers, qui charge le modèle et son
découpeur en jetons.

\subsection{La tâche et la récompense}

La tâche retenue, \file{mul2}, demande le produit de deux entiers de 12 à 99. Le fichier énumère les
$88 \times 88 = 7\,744$ couples ordonnés, les mélange avec un générateur de graine aléatoire fixe
(\texttt{random.Random(12345)}), indépendant de la graine aléatoire de l'essai, puis réserve les 300 premiers
au test. Les 7\,444 autres servent à l'entraînement. L'énoncé suit un gabarit fixe, \enquote{What is
$a$ * $b$? Think step by step, then finish with the final answer as an integer.}, passé dans le
gabarit de conversation du modèle.

Une réponse est juste quand le dernier entier qu'elle contient, virgules de milliers retirées, est
égal au produit. Aucun format n'est exigé, pour ne pas récompenser le format. Soit $c(o, q) \in
\{0,1\}$ ce verdict pour la réponse $o$ au problème $q$. Les deux conditions du protocole
définissent la récompense $r$ de chaque réponse :
\begin{align}
  \text{vraie récompense (\texttt{true})} &: \quad r = c(o, q),\\
  \text{récompense tirée au sort (\texttt{random})} &: \quad r \sim \mathrm{Bernoulli}(1/2),
  \text{ tirée sans regarder } o.
\end{align}
Le tirage de la pièce et celui des problèmes de chaque pas viennent du même générateur,
\texttt{random.Random(graine)}, qui dépend de la seule graine aléatoire de l'essai.

\subsection{Un pas d'entraînement}

Chaque pas tire $B = 4$ problèmes d'entraînement, sans remise à l'intérieur du pas : un problème peut revenir
d'un pas à l'autre. Pour chaque problème $q_b$, la
politique courante $\pi_{\theta_{\text{old}}}$ produit $G = 8$ réponses $o_{b,1}, \dots, o_{b,8}$,
échantillonnées à la température 1, sans \emph{top-p} ni \emph{top-k}, limitées à 160 jetons. Pour Llama, la
distribution échantillonnée est celle du modèle. Pour Qwen, elle ne l'est pas (écart~18, point~a). Le fichier
\file{generation\_config.json} de Qwen2.5 fixe une pénalité de répétition de 1,1, que \file{grpo.py} ne
remplace pas. Avant chaque tirage, le logit $\ell$ de tout jeton déjà présent dans le texte devient $\ell/1{,}1$
s'il est positif et $1{,}1\,\ell$ sinon. Trois conséquences en découlent. Les réponses de Qwen ne suivent pas
la politique dont $\log \pi_{\theta_{\text{old}}}$ donne les probabilités, sans correction. La pénalité frappe
les chiffres répétés d'un calcul. La comparaison entre les deux familles mêle la famille et ce réglage de
décodage. Chaque réponse reçoit sa récompense $r_{b,i}$. L'avantage
d'une réponse compare sa récompense à celles de son groupe :
\begin{equation}
  \hat A_{b,i} = \frac{r_{b,i} - \bar r_b}{s_b + 10^{-4}}, \qquad
  \bar r_b = \frac1G \sum_{j=1}^{G} r_{b,j}, \qquad
  s_b = \Bigl(\frac{1}{G-1}\sum_{j=1}^{G} (r_{b,j} - \bar r_b)^2\Bigr)^{1/2}.
\end{equation}
L'écart type $s_b$ suit la convention de PyTorch, avec le facteur $G-1$. Un groupe où les huit
réponses reçoivent la même récompense donne un avantage nul à chacune et n'apporte aucun gradient.

Pour le jeton $t$ de la réponse $o_{b,i}$, le rapport de probabilités compare la politique en cours
de mise à jour à celle qui a produit la réponse :
\begin{equation}
  \rho_{b,i,t}(\theta) = \exp\bigl(\log \pi_\theta(o_{b,i,t} \mid q_b, o_{b,i,<t})
    - \log \pi_{\theta_{\text{old}}}(o_{b,i,t} \mid q_b, o_{b,i,<t})\bigr).
\end{equation}
La perte minimisée est l'opposé de l'objectif coupé, moyenné par jeton dans chaque réponse puis
par réponse dans le pas :
\begin{equation}
  \mathcal L(\theta) = -\frac{1}{BG} \sum_{b=1}^{B} \sum_{i=1}^{G} \frac{1}{|o_{b,i}|}
  \sum_{t=1}^{|o_{b,i}|} \min\Bigl(\rho_{b,i,t}(\theta)\,\hat A_{b,i},\;
  \operatorname{clip}\bigl(\rho_{b,i,t}(\theta), 1-\varepsilon, 1+\varepsilon\bigr)\,\hat A_{b,i}\Bigr),
  \qquad \varepsilon = 0{,}2.
  \label{eq:grpo}
\end{equation}
La longueur $|o_{b,i}|$ compte les jetons de la réponse qui ne sont pas du remplissage (le masque
\texttt{target != pad\_token\_id} de la fonction \texttt{token\_logps}). Chez Qwen, le jeton de remplissage
est aussi le jeton d'arrêt \texttt{<|endoftext|>} : le masque retire donc du gradient le jeton qui termine la
réponse (écart~18, point~d). Les log-probabilités de
référence $\log \pi_{\theta_{\text{old}}}$ sont calculées une fois, sans gradient, avant la première
mise à jour.

Chaque lot de 32 réponses sert à deux passes de mise à jour (\texttt{EPOCHS = 2}). À chaque passe,
le gradient s'accumule sur des sous-lots de huit réponses (deux pour Llama, réglage \texttt{MICRO}),
ce qui donne le même gradient avec moins de mémoire. La norme du gradient est ensuite coupée à 1,
puis l'optimiseur Adam \parencite{kingma2015adam}, avec les réglages par défaut de PyTorch
($\beta_1 = 0{,}9$, $\beta_2 = 0{,}999$, $\varepsilon = 10^{-8}$), applique un taux d'apprentissage de
$3 \times 10^{-6}$. Les poids restent en virgule flottante sur 32 bits.

\begin{notebox}
\textbf{Ce qui distingue notre GRPO de celui de \textcite{shao2024deepseekmath}.} Notre perte n'a
aucun terme de divergence de Kullback-Leibler par rapport au modèle de départ, ce que le protocole
du 8 novembre 2026 annonçait. Le groupe compte 8 réponses par problème et le pas 4 problèmes, bien
moins que dans les articles de référence. La limite de 160 jetons est très courte pour un
raisonnement écrit. Le taux d'apprentissage n'a pas été choisi par le pilote prévu, interrompu au
pas 10 (écarts 4 et 6). Ces choix suffisent à une illustration. Ils interdisent de lire nos chiffres
comme une mesure de GRPO à l'échelle des articles.
\end{notebox}

\subsection{Les garde-fous contre un calcul corrompu}

Les premiers essais sur le processeur graphique du Mac (GPU, via l'interface MPS de PyTorch) ont
planté sur des probabilités non finies. L'écart~10 du protocole en documente la cause : macOS
interrompt un calcul GPU qui gêne l'usage interactif de la machine. Ce calcul revient parfois
sans erreur, avec des valeurs corrompues. Le même pas, refait sur le processeur central (CPU), donne
une norme de gradient de 14,4 avant écrêtage, sans aucune valeur non finie. Depuis cet écart,
\file{grpo.py} applique quatre règles :
\begin{enumerate}
  \item une mise à jour dont la norme du gradient avant écrêtage est non finie ou dépasse 1\,000 est
        recalculée, jusqu'à cinq fois de suite, puis l'essai s'arrête (\texttt{GRAD\_SANE},
        \texttt{RETRIES}) ;
  \item une génération qui lève une erreur GPU est refaite, jusqu'à cinq fois ;
  \item l'essai s'arrête si les log-probabilités de référence ou les poids deviennent non finis ;
  \item le journal compte les recalculs à chaque pas (champ \texttt{retries}).
\end{enumerate}
Le seuil de 1\,000 vient de la norme de 14,4 mesurée sur CPU. Un gradient corrompu mais fini, sous
ce seuil, puis écrêté à 1, passerait sans trace : le protocole le note comme une limite, en
particulier pour les trois premiers essais, qui ont tourné avant ces règles.

\subsection{La mesure}

L'évaluation se fait sur le jeu de test et alterne deux modes. La justesse gloutonne (\emph{greedy}) prend une seule
réponse sans hasard, le jeton le plus probable à chaque pas, pour chacun des 150 premiers problèmes
de test. La justesse échantillonnée tire
8 réponses à la température 1 pour chacun des 50 premiers problèmes de test. Avec $c_{jk}$ le
verdict de la réponse $k$ au problème $j$ et $\tau_{jk}$ l'indicateur d'une réponse coupée :
\begin{equation}
  \text{pass@1} = \frac{1}{50}\sum_{j=1}^{50} \frac18 \sum_{k=1}^{8} c_{jk}, \qquad
  \text{pass@8} = \frac{1}{50}\sum_{j=1}^{50} \max_{k} c_{jk}, \qquad
  \text{part coupée} = \frac{1}{400}\sum_{j,k} \tau_{jk}.
\end{equation}
Une réponse est coupée quand ses 160 jetons ne contiennent ni jeton de fin ni jeton de remplissage
(fonction \texttt{truncated}). Cette fonction ne connaît que le jeton de fin principal : elle ignore les deux
jetons d'arrêt supplémentaires de Llama, 128001 et 128008 (écart~18, point~d). Chez Qwen, la pénalité de
répétition s'applique aussi à la réponse gloutonne. Le pass@8 est ici la fréquence observée sur exactement huit tirages,
sans l'estimateur sans biais employé quand on tire plus de réponses que $k$. L'évaluation a lieu au
pas 0, au pas 50 et au pas 100 (\texttt{EVAL\_EVERY = 50}). Tous les dix pas, le journal note aussi
la part de réponses justes et coupées du lot d'entraînement, la perte, la norme du gradient, les
recalculs et le temps écoulé. Chaque essai écrit un fichier JSONL,
\file{sim/llm/results/<modèle>/mul2/<récompense>/seed<graine>.jsonl}. Le journal garde les quatre mesures
de chaque évaluation, mais aucun verdict par problème. L'évaluation tire ses réponses dans le même générateur
aléatoire que l'entraînement : l'évaluation du pas 50 change les tirages des pas suivants (écart~18,
point~d).

La mesure principale est le pass@1 échantillonné, celle que GRPO entraîne (écart 3). La première version
de ce chapitre laissait la justesse gloutonne et le pass@8 dans les journaux. L'écart~19 les rapporte, avec
une cinquième mesure, la justesse parmi les réponses finies :
\begin{equation}
  \text{justesse des finies} = \frac{\text{pass@1}}{1 - \text{part coupée}},
\end{equation}
qui suppose qu'une réponse coupée n'est jamais juste (0 sur 23 dans le diagnostic de l'écart~5,
\cref{sec:llm-coupure}). Elle sépare ce que pass@1 doit à la justesse de ce qu'il doit à la brièveté. Le
pass@8 répond à la question de \textcite{yue2025rlvr} : l'affinage élargit-il ce que le modèle sait faire, ou
concentre-t-il sa probabilité sur des réponses qu'il trouvait déjà en plusieurs tirages ?

\section{Le protocole et ses essais}
\label{sec:llm-protocole}

\paragraph{Modèles} Deux familles, deux petits modèles ouverts : Qwen2.5 de 0,5 milliard de
paramètres, version \emph{Instruct} \parencite{qwen2025qwen25}, puis Llama 3.2 de 1 milliard,
version \emph{Instruct}, dans la copie publiée par \texttt{unsloth}
\parencite{meta2024llama32,grattafiori2024llama3}. Depuis l'écart 15 (c), \file{grpo.py} fixe la
révision exacte des fichiers téléchargés (\texttt{7ae5576\dots} pour Qwen, \texttt{5a8abab\dots}
pour Llama), les seules présentes dans le cache de la machine, téléchargées les 22 et 27 octobre
2026, avant le premier essai.

\paragraph{Pilote} Un pilote sans entraînement devait fixer la tâche : celle où les deux modèles se situent
entre 15~\% et 70~\% de justesse. Le protocole le voulait sur des problèmes hors test. Il a tourné sur le test
(écart~18, point~b) : l'option \texttt{-{}-pilot} appelle la fonction d'évaluation, qui lit les 150 et 50
premiers problèmes du test. Le choix de \file{mul2} et celui de la limite de 160 jetons ont donc été faits en
regardant le test. Aucun entraînement n'a vu le test. Aucune tâche ne place les deux
modèles dans cette plage en glouton. La somme de trois nombres (\file{add3}) est trop facile pour
Llama. La multiplication \file{mul2} donne en glouton 7~\% (Qwen) et 73~\% (Llama), en pass@1
échantillonné 11~\% et 23~\% (écart 3). Le même pilote a montré que 64 jetons coupaient les réponses de
Qwen, qui en écrit 100 à 160, d'où 0,7~\% de justesse apparente : la limite est passée à 160 jetons,
au vu de vraies sorties, avant tout entraînement (écart 1).

\paragraph{Plan} Deux modèles, deux récompenses, 100 pas par essai, taux d'apprentissage $3 \times
10^{-6}$ pour toutes les conditions. Le protocole prévoyait 200 pas et cinq graines aléatoires.
L'écart 6 a réduit le plan à 100 pas et trois graines aléatoires (0 à 2), avant tout essai retenu, pour tenir
dans le calcul disponible. L'écart 14, écrit le 15 novembre 2026 avant tout nouvel essai, a rétabli
les graines aléatoires 3 et 4, avec le même code et les mêmes réglages. Le point de départ de chaque essai est
le modèle de base, mesuré avec la graine aléatoire de l'essai.

\paragraph{Statistique} Pour chaque graine aléatoire $s$, le gain vaut $g_s = \text{pass@1}_{100} -
\text{pass@1}_{0}$. Chaque configuration (modèle, récompense) se résume par la moyenne des gains et
un intervalle bootstrap par percentile à 95~\% : 10\,000 rééchantillonnages avec remise des gains,
générateur \texttt{random.Random(0)}, bornes aux rangs 250 et 9\,750 des moyennes triées
(\file{sim/llm/summary.py}). Les détails du bootstrap ont été fixés après les premiers résultats,
sans choix entre plusieurs variantes (écart 13). L'intervalle est très grossier : trois graines aléatoires
n'offrent que dix rééchantillons distincts, cinq graines aléatoires 126.

\paragraph{Statistique ajoutée} L'écart~19, écrit avant le script \file{sim/llm/logs\_stats.py}, ajoute
trois outils pour les cinq mesures. Une partie de ces valeurs avait été affichée par la relecture technique
du même jour : la règle de lecture ne dépend pas de leur signe.
\begin{itemize}
  \item Dans chaque configuration, à côté de l'intervalle par percentile, l'intervalle de Student
    $\bar g \pm t_{n-1;\,0{,}975}\, s_g/\sqrt n$, où $\bar g$ et $s_g$ sont la moyenne et l'écart type des
    $n$ gains. Avec cinq graines aléatoires, l'intervalle par percentile n'a qu'environ 60~\% de la largeur de
    l'intervalle de Student et couvre trop peu. Le test exact du signe compare les gains aux $2^n$ jeux de signes possibles.
  \item Entre les deux récompenses, la question même du protocole : la différence des gains moyens, son
    intervalle de Welch, avec les degrés de liberté de Welch-Satterthwaite, puis un test exact de permutation
    sur toutes les répartitions des gains entre les deux récompenses, 252 pour Qwen et 126 pour Llama. La
    correction de Holm s'applique aux deux modèles pour pass@1, le contraste principal, comme pour la justesse
    des réponses finies. Elle s'applique ensuite aux huit autres contrastes.
  \item Une règle : un gain ou un contraste est \emph{établi} quand son intervalle de Student ou de Welch
    exclut 0 et que son $p$ exact, corrigé par Holm, passe sous 0,05.
\end{itemize}
Cette règle a une conséquence que le protocole n'avait pas vue. Avec cinq graines aléatoires, le test exact du
signe ne descend jamais sous $p = 2/2^5 = 0{,}0625$, atteint quand les cinq gains ont le même signe. Avec
quatre, il ne descend pas sous $2/2^4 = 0{,}125$. Aucun gain à l'intérieur d'une configuration ne peut donc
être établi avec cinq graines aléatoires, même quand toutes vont dans le même sens. Plus de réponses par
problème n'y changerait rien. Au pas 0, l'écart type de pass@1 entre graines aléatoires vaut 0,7~point pour
Qwen et 1,1~point pour Llama, sous l'erreur type binomiale de 400 réponses (1,5 et 2,1~points). Les gains
sous vraie récompense ont un écart type d'environ 24~points pour Qwen et 10 pour Llama. Il faut plus de
graines aléatoires.

\paragraph{Un essai qui échoue deux fois} L'essai Llama à récompense tirée au sort et graine aléatoire 4
s'est arrêté deux nuits de suite sur le premier garde-fou, dès sa première mise à jour. L'écart 16
compte ce second échec comme un résultat : la configuration se résume sur ses quatre essais finis,
l'échec écrit à côté. Une analyse de sensibilité remplace aussi le gain manquant par le plus petit
puis par le plus grand gain des dix-neuf essais finis.

\paragraph{Matériel} Tous les essais ont tourné sur le Mac du projet
(\cref{chap:08-logiciel,chap:D-environnements}). Les trois premiers essais Qwen ont tourné sur le
GPU, tôt le matin, machine au repos. Après les plantages de l'écart 8 et le diagnostic de l'écart 10, les essais Qwen
restants sont passés sur CPU, avec deux fils d'exécution (\file{cpu.sh}, écart 11). Les essais
Llama, environ deux fois et demie plus lents que Qwen sur CPU, sont restés sur GPU, la nuit
seulement (\file{night.sh}), avec les garde-fous. Le \cref{tab:grpo} donne l'appareil de chaque
essai.

\section{Les résultats}
\label{sec:llm-resultats}

\gentable{grpo}

Le \cref{tab:grpo} donne, pour chacun des vingt essais, le pass@1 au départ et après 100 pas,
l'appareil et la part de réponses coupées à la fin. La \cref{fig:grpo-graines} montre chaque essai
par une flèche, la \cref{fig:grpo-courbes} les trois points d'évaluation de chaque essai.

\fig{grpo-graines}{Pass@1 au départ et après cent pas d'entraînement, pour chaque modèle, chaque
récompense et chacune des cinq graines aléatoires. Un rond plein désigne un essai sur GPU, un carré
vide un essai sur CPU, une croix l'essai arrêté deux fois. Avec la vraie récompense, le même
programme donne des résultats très différents d'une graine aléatoire à l'autre. À pile ou face, la
réussite tombe à zéro ou presque.}

\fig{grpo-courbes}{Pass@1 de test aux pas 0, 50 et 100 de chaque essai, par modèle, récompense et
graine aléatoire. L'évaluation n'a lieu que tous les cinquante pas : la courbe relie trois points.}

\subsection{La récompense tirée au sort}

À pile ou face, la réussite tombe à 1,25~\% ou moins dans les neuf essais finis. Le gain moyen vaut
$-10$ points sur Qwen (de $-11$ à $-9$) et $-24$ points sur Llama (de $-24$ à $-23$). Remplacer le
gain manquant de Llama par le plus petit gain des dix-neuf essais finis change peu ($-24$, de $-24$ à
$-23$). Le remplacer par le plus grand, $+51$, relève la moyenne à $-9$, avec un intervalle de $-24$ à
$+21$. Chez nous, à cette échelle, la récompense tirée au sort détruit la compétence de départ des
deux familles. Elle ne reproduit pas le gain que \textcite{shao2025spurious} observent sur
Qwen2.5-Math-7B. Nos conditions diffèrent des leurs sur presque tout (taille et spécialisation du
modèle, tâche, longueur des réponses, nombre de pas, absence de terme KL) : ce désaccord ne réfute
pas leur résultat.

\subsection{La vraie récompense, en deux temps}

Avec trois graines aléatoires, Qwen part d'environ 10~\% et finit à 36~\%, 61~\% et 13~\%. Le gain
moyen vaut $+26$ points, avec un intervalle de $+2$ à $+51$. Sur Llama, la réussite part de 23~\% à
25~\% et finit à 20~\%, 48~\% et 25~\%, soit un gain moyen de $+7$ points, de $-4$ à $+22$. Le même
programme donne trois résultats très différents : c'est la faille F11, même code, autre graine
aléatoire, en petit.

Avec deux graines aléatoires de plus, les nouveaux essais Qwen finissent à 13~\% et 0~\%. Sur cinq
graines aléatoires, le gain moyen tombe à $+14$ points, de $-3$ à $+34$ : l'intervalle contient
désormais zéro. Le gain que nous lisions sur Qwen avec trois graines aléatoires ne tient plus. Sur Llama, les
deux nouveaux essais finissent à 27~\% et 33~\%. Le gain moyen vaut $+6$ points, de $-0{,}4$ à $+15$ :
zéro reste dans l'intervalle. Les intervalles de Student de l'écart~19 sont plus larges : de $-16$ à $+45$
pour Qwen, de $-6$ à $+19$ pour Llama, avec un $p$ de signe de 0,31 pour les deux. Avec cinq graines
aléatoires, nous ne pouvons pas dire que GRPO fait progresser l'un ou l'autre modèle sur cette tâche, dans
ces conditions.

\subsection{Les mesures que le pass@1 cachait}

Le \cref{tab:grpo-mesures} donne les cinq mesures de chaque configuration, la \cref{fig:grpo-mesures} leurs
gains graine par graine. Sous la vraie récompense, Llama gagne 6~points de pass@1 pendant que sa justesse
gloutonne chute dans ses cinq essais, de 72,1~\% à 29,1~\% en moyenne ($-43$~points, Student
$[-59\,;\,-27]$). Son pass@8 chute aussi, de 75,2~\% à 44,8~\% ($-30$~points, $[-44\,;\,-17]$). Le modèle
ne calcule pas mieux. La distribution qu'il échantillonne à la température 1 se resserre : un tirage au
hasard tombe un peu plus souvent juste, alors que la réponse la plus probable et la meilleure de huit
tombent moins souvent juste. Chez Qwen, le pass@8 baisse dans quatre essais sur cinq, de 51,6~\% à 34,4~\%,
avec un intervalle de Student qui contient zéro. À la question de \textcite{yue2025rlvr}, nos essais
répondent que le pass@8 ne monte en moyenne dans aucune configuration. Sous la règle de l'écart~19, ces chutes restent non établies, faute
d'un $p$ de signe sous 0,0625, alors que leurs intervalles de Student excluent zéro. Les deux constats valent
ensemble.

\gentable{grpo-mesures}

\fig{grpo-mesures}{Gains par graine aléatoire en trois mesures, entre le pas 0 et le pas 100, pour Qwen2.5 de
0,5 milliard de paramètres et Llama 3.2 de 1 milliard, sous la vraie récompense et sous la récompense tirée à
pile ou face. Un point par graine aléatoire, la moyenne, l'intervalle de Student (trait épais) et
l'intervalle par percentile du protocole (trait fin). Sous la vraie récompense, le pass@1 de Llama monte de
6~points pendant que sa justesse gloutonne perd 43~points et son pass@8 30 : la distribution échantillonnée
se resserre, le modèle ne multiplie pas mieux. Avec cinq graines aléatoires, l'intervalle par percentile est
bien plus étroit que celui de Student.}

\subsection{Vraie récompense contre pile ou face}

Le protocole posait une question que la première version de ce chapitre ne testait pas : la vraie récompense
fait-elle mieux que la récompense tirée au sort ? Le \cref{tab:grpo-contraste} y répond par la différence
des gains moyens. Pour Llama, la vraie récompense fait mieux en pass@1, de $+30$~points
$[+18\,;\,+43]$ (permutation exacte $p = 0{,}008$, Holm 0,016) : c'est établi. Pour Qwen, l'écart va dans le
même sens, $+25$~points $[-6\,;\,+55]$ ($p = 0{,}040$, Holm 0,040), mais son intervalle contient zéro : il
n'est pas établi. Aucun des huit contrastes secondaires ne passe Holm. La raison du contraste de Llama
compte autant que son signe. La vraie récompense ne fait pas progresser Llama : elle le dégrade moins que la
récompense tirée au sort, qui le détruit (pass@1 de 24~\% à 0,3~\%, justesse gloutonne de 72~\% à 0,8~\%).

\gentable{grpo-contraste}

\subsection{Le matériel et la graine aléatoire, mêlés}
\label{sec:llm-materiel}

Le gain de Qwen avec la vraie récompense vient des deux essais faits sur GPU, graines aléatoires 0
et 1 : $+26$ et $+51$ points. Les trois essais à vraie récompense qui ont tourné sur CPU (graines aléatoires 2
à 4) ont un gain moyen de $-2$ points, de $-11$ à $+4$. Les quatre essais sur CPU à récompense tirée
au sort (graines aléatoires 1 à 4) ont un gain moyen de $-11$ points, de $-11$ à $-10$. Sur CPU seul, la vraie
récompense n'apporte à Qwen aucun gain détectable. Données :
\file{sim/llm/results/gains\_5.json} et \file{extra\_5.json}.

Ces chiffres ne séparent pas deux causes. Le matériel a changé avec la graine aléatoire : chaque
graine aléatoire de Qwen à vraie récompense n'a tourné que sur un appareil. Deux explications restent
ouvertes. Selon la première, le hasard des graines aléatoires suffit : les deux graines favorables sont tombées
sur le GPU. Selon la seconde, l'appareil change la trajectoire de l'entraînement, par l'arrondi des
calculs ou par des calculs GPU corrompus mais finis, que les premiers essais n'auraient pas
détectés. Cinq essais ne départagent pas ces explications.

Les journaux établissent deux faits sur l'appareil (écart~18, point~c). Le pass@1 au pas 0 dépend
du modèle de base, de la graine aléatoire et du hasard des tirages, puisque la récompense n'intervient qu'après
cette première évaluation. La justesse gloutonne au pas 0, qui ne tire rien au hasard, vaut \num{0.0733} dans
les dix essais Qwen, sur CPU comme sur GPU : l'appareil ne change pas la réponse la plus probable du modèle de
départ. Les écarts de pass@1 au pas 0 entre essais Qwen d'une même graine aléatoire viennent des flux de
hasard, qui diffèrent entre CPU et GPU pour une même graine. Pour Qwen graine aléatoire 1, l'essai à vraie
récompense, sur GPU, part de \num{0.1025}, l'essai à récompense tirée au sort, sur CPU, de \num{0.1175}, pour
cette raison. Le signe direct du non-déterminisme du calcul MPS vient de Llama. Pour la graine aléatoire 3,
les deux essais ont tourné sur GPU, avec le même modèle et le même hasard. Ils partent pourtant d'un pass@1 de
\num{0.25} et de \num{0.2425}. Leur justesse gloutonne vaut \num{0.707} d'un côté, \num{0.667} de l'autre. L'écart~7
l'avait déjà constaté (\cref{tab:grpo}).

Les journaux montrent aussi des groupes sans gradient. Sous la vraie récompense, beaucoup de groupes ont huit
réponses toutes fausses, ou toutes justes, donc un avantage nul pour chacune. La norme du gradient vaut zéro
dans 5, 4 et 7 des 10 pas journalisés pour les essais Qwen des graines aléatoires 2, 3 et 4 (les graines
aléatoires 0 et 1 ne l'ont pas journalisée). Elle vaut zéro dans 3, 0, 3, 3 et 2 pas sur 10 pour les cinq
essais Llama. L'essai Qwen de la graine
aléatoire 4 cesse d'apprendre dès le pas 50, à 0~\% de pass@1, de pass@8 et de justesse gloutonne. Sous la
récompense tirée au sort, la pièce donne presque toujours des récompenses mêlées : aucun pas journalisé n'a
un gradient nul.

\paragraph{Un essai hors protocole} L'écart 16 (e) prévoyait de lancer une fois sur CPU l'essai
Llama à récompense tirée au sort et graine aléatoire 4, après tous les autres calculs, pour savoir si son
double échec tenait au GPU. L'essai a abouti le 21 novembre 2026, sans aucun recalcul. Son pass@1
passe de \num{0.2325} au pas 0 à \num{0.0025} au pas 100
(\file{sim/llm/results/llama1\_cpu/mul2/random/seed4.jsonl}). Le résultat sort de toute
statistique, comme l'écart l'avait fixé. Il cadre avec un échec dû au GPU, sans le prouver : un
seul essai ne sépare pas l'appareil du hasard. Sur CPU, ce dixième essai se comporte comme les neuf autres essais à récompense tirée au sort : la réussite s'effondre.

\section{Le piège de la coupure}
\label{sec:llm-coupure}

\subsection{Le diagnostic}

L'incident de l'écart 4 a mis ce piège au jour. Pendant le pilote du taux d'apprentissage, la part de
réponses justes du lot d'entraînement valait 0,91 au pas 10, alors que le pass@1 de test du modèle
de base valait 0,11. Le diagnostic de l'écart 5 (\file{sim/llm/diag/replay\_batch.py}) rejoue,
sans entraînement, le lot exact du pas 10 du pilote : les quatre multiplications $48 \times 49$,
$14 \times 36$, $21 \times 50$ et $45 \times 17$, huit réponses chacune, avec le modèle Qwen de
départ. Le modèle donne 5 réponses justes sur 32, ce qui concorde avec le test. Mais 23 réponses sur 32
atteignent la limite de 160 jetons avant de donner un résultat. Aucune de ces réponses coupées n'est
juste. Parmi les 9 réponses complètes, 5 sont justes (\cref{fig:grpo-coupure,fig:grpo-truncation}).

\fig{grpo-coupure}{Les réponses du modèle Qwen2.5 de 0,5 milliard de paramètres avant tout
entraînement, huit par problème, sur les quatre multiplications du dixième pas du pilote. Une réponse
complète et juste est en vert, une réponse complète et fausse en gris, une réponse coupée en rouge.
La plupart atteignent la limite de 160 jetons avant de donner un résultat et comptent pour fausses.}

\fig[0.8\linewidth]{grpo-truncation}{Les mêmes trente-deux réponses, en barres empilées : pour chaque
problème, la part des huit réponses complètes et justes, complètes et fausses, coupées à 160 jetons.}

Le 0,91 du pas 10 n'était pas une erreur de décompte. L'explication la plus simple : la vraie
récompense apprend d'abord au modèle à terminer avant la coupure, si bien que le lot d'entraînement
était devenu plus facile à finir. La récompense tirée au sort n'a pas ce levier, puisqu'elle ne voit pas la réponse. Le protocole a gardé
les 160 jetons et ajouté la mesure de la part coupée, à l'entraînement et au test, avant tout essai
retenu.

\subsection{Ce que montrent les vingt essais}

Au départ, la part de réponses coupées est élevée pour Qwen et faible pour Llama : de 71~\% à
77~\% pour Qwen, 3~\% au plus pour Llama, d'après les journaux d'évaluation au pas 0. Avec la vraie récompense, sur les dix essais, elle finit entre 0~\% et 29~\%. À
pile ou face, elle finit au-dessus de la moitié dans six essais sur neuf, entre 2~\% et 16~\% dans
les trois autres. Deux essais Llama à récompense tirée au sort finissent avec 99~\% et 100~\% de
réponses coupées. Un troisième n'en coupe que 2~\% et ne donne plus aucune bonne réponse : la
récompense tirée au sort ne se contente pas d'allonger les réponses au-delà de la limite.

Une réponse coupée compte pour fausse. Un modèle qui apprend à répondre plus court gagne donc des
points de pass@1 sans mieux calculer. La première version de ce chapitre ne pouvait pas chiffrer cette part.
La justesse parmi les réponses finies de l'écart~19 la chiffre (\cref{tab:grpo-mesures}).
\begin{itemize}
  \item Chez Qwen, sous la vraie récompense, la part coupée tombe de 74,9~\% à 6,0~\% dans les cinq essais.
    La justesse parmi les réponses finies \emph{baisse}, de 41,2~\% à 27,5~\% ($-14$~points, Student
    $[-46\,;\,+19]$, deux essais sur cinq en hausse). Tout le gain de pass@1 de Qwen s'explique par la fin des
    réponses coupées. Le modèle apprend à finir avant 160 jetons. Il ne multiplie pas mieux.
  \item Chez Llama, la longueur ne bouge presque pas : la part coupée passe de 2,0~\% à 2,7~\%. La justesse
    des réponses finies suit le pass@1, $+7$~points, Student $[-5\,;\,+19]$.
\end{itemize}
À l'intérieur de la configuration Qwen à vraie récompense, raccourcir plus ne donne pas plus de gain : la
corrélation de Pearson entre la variation de la part coupée et celle de pass@1 vaut $+0{,}23$ sur cinq
essais. Sur les deux récompenses ensemble, la corrélation de Spearman vaut $-0{,}63$ pour Qwen et $-0{,}83$
pour Llama, ce qui reflète surtout l'effet de la récompense (\file{sim/llm/results/logs\_stats\_5.json}, clé
\texttt{length}).

\begin{keybox}
\textbf{Leçon pour un banc d'essai de modèle de langue.} Une limite de longueur transforme une mesure
de justesse en une mesure mixte de justesse et de concision. Un article qui rapporte un gain de
pass@1 se lit avec la limite de longueur, la part de réponses coupées avant et après
l'entraînement et la justesse des réponses complètes. Sans ces trois nombres, le lecteur ignore ce
que le gain mesure.
\end{keybox}

\section{Ce que ce volet établit et ce qu'il n'établit pas}
\label{sec:llm-portee}

Ce volet établit quatre faits, pour deux petits modèles, une tâche de multiplication et 100 pas.
\begin{itemize}
  \item Avec la vraie récompense, l'écart entre graines aléatoires est aussi grand que le gain moyen de
    pass@1, comme dans le RL classique de 2018. Aucun gain de pass@1 n'est établi.
  \item Une récompense tirée au sort fait chuter la réussite des deux familles, dans tous les essais et dans
    les quatre mesures de justesse. C'est elle qui donne à Llama son seul contraste établi entre les deux récompenses.
  \item Chez Llama, la vraie récompense resserre la distribution échantillonnée : le pass@1 monte un peu, la
    justesse gloutonne et le pass@8 chutent dans les cinq essais.
  \item Chez Qwen, tout le gain de pass@1 est un gain de brièveté : la justesse des réponses finies ne monte
    pas.
\end{itemize}
Deux défauts déclarés pèsent sur ces faits : la pénalité de répétition de Qwen et un pilote qui a lu le test
(écart~18). Le protocole à cinq graines aléatoires ne peut établir aucun gain dans une
configuration. Une relance est écrite d'avance pour lever ces défauts (écart~20) : pénalité de répétition à
1, pilote tenu hors du test, dix graines aléatoires par configuration sur le seul GPU, verdicts par problème
gardés et une réponse gloutonne à 512 jetons, que la limite de longueur ne touche pas. Lancée le 22 novembre
2026, elle attend que le GPU se libère et n'a encore produit aucun résultat. Sa règle est fixée : si le
contraste principal n'est établi pour aucun des deux modèles, ce volet passe en annexe.

Il n'établit ni l'effet de GRPO à l'échelle des articles, ni l'absence de l'effet de
\textcite{shao2025spurious} sur des modèles plus grands ou spécialisés en mathématiques. Il ne
sépare pas l'effet du matériel de celui de la graine aléatoire pour Qwen. Les limites qui en
découlent sont reprises au \cref{chap:10-limites}.
